\begin{ouverture}[courbe]{2}{Puissances, racines et écriture scientifique}{Partie 1 · Nombres et calculs}
{Exposants · Inverses · Ordres de grandeur · Racines carrées}
\objectif{Utiliser les puissances entières.}
\objectif{Écrire et calculer en notation scientifique.}
\objectif{Résoudre un problème avec une racine carrée.}
\prerequis{Produits, quotients et carrés.}
\end{ouverture}
\pagestyle{otMain}\essentiel\fontsize{11}{15}\selectfont
\essSection{1}{Puissances et ordres de grandeur}
\begin{definition}Pour $n$ entier positif, $a^n$ est un produit de $n$ facteurs égaux à $a$. Pour $a\ne0$, $a^0=1$ et $a^{-n}=1/a^n$.\end{definition}
\begin{propriete}Pour les bases et dénominateurs autorisés :
\[a^ma^n=a^{m+n},\quad \frac{a^m}{a^n}=a^{m-n},\quad
(a^m)^n=a^{mn},\quad (ab)^n=a^nb^n.\]\end{propriete}
\begin{definition}La notation scientifique d’un nombre positif est $c\times10^p$, avec $1\le c<10$ et $p$ entier. Par exemple $0{,}000056=5{,}6\times10^{-5}$.\end{definition}
\begin{savoirfaire}[Calculer en notation scientifique]Regrouper les coefficients, puis les puissances de 10. $(3\times10^5)(4\times10^{-2})=12\times10^3=1{,}2\times10^4$.\end{savoirfaire}
\begin{attention}$a^m+a^n$ ne se transforme pas en $a^{m+n}$. Les règles sur les exposants concernent des produits, quotients et puissances.\end{attention}
\clearpage
\essSection{2}{Racines carrées et mesures}
\begin{definition}Pour $a\ge0$, $\sqrt a$ est l’unique nombre positif ou nul dont le carré vaut $a$. Ainsi $(\sqrt a)^2=a$. Une racine d’un nombre négatif n’est pas un nombre réel.\end{definition}
\begin{exemple}Un carré d’aire 45 cm$^2$ a pour côté $\sqrt{45}$ cm, environ $6{,}71$ cm. La longueur est positive ; l’équation $x^2=45$ a aussi la solution négative $-\sqrt{45}$.\end{exemple}
\begin{propriete}Pour $a,b\ge0$, $\sqrt{ab}=\sqrt a\sqrt b$. Par exemple $\sqrt{72}=6\sqrt2$. En revanche, $\sqrt{a+b}$ n’est pas généralement $\sqrt a+\sqrt b$.\end{propriete}
\begin{savoirfaire}[Arrondir une mesure]Conserver les valeurs exactes pendant les étapes. Arrondir à la précision demandée à la fin, dans l’unité choisie. Un encadrement comme $6{,}7<\sqrt{45}<6{,}8$ se vérifie en comparant les carrés positifs.\end{savoirfaire}
\begin{attention}$\sqrt{(-4)^2}=4$. La racine carrée renvoie toujours un nombre positif ou nul. Pour $(-4)^2$, les parenthèses font porter le carré sur le signe.\end{attention}
\clearpage\rubrique{Découverte}
\begin{blocLivre}{ex-2-1}\begin{exercice}[1]{Comprendre un exposant}
Écrire $7\times7\times7\times7$ comme une puissance. Calculer $2^{-3}$ et $9^0$.
\end{exercice}\end{blocLivre}
\begin{blocLivre}{ex-2-2}\begin{exercice}[1]{Parenthèses et signe}
Calculer $(-5)^2$, $-5^2$, $(-2)^3$ et $(-2)^4$.
\end{exercice}\end{blocLivre}
\begin{blocLivre}{ex-2-3}\begin{exercice}[1]{Règles de puissances}
Simplifier $3^4\times3^2$, $5^7/5^3$ et $(2^3)^4$.
\end{exercice}\end{blocLivre}
\begin{blocLivre}{ex-2-4}\begin{exercice}[1]{Écriture scientifique}
Écrire $74000000$ et $0{,}000092$ en notation scientifique.
\end{exercice}\end{blocLivre}
\begin{blocLivre}{ex-2-5}\begin{exercice}[1]{Quelques racines}
Calculer $\sqrt{144}$, $\sqrt{0{,}81}$ et $\sqrt{49/64}$.
\end{exercice}\end{blocLivre}

\clearpage\rubrique{Entraînement}
\begin{blocLivre}{ex-2-6}\begin{exercice}[2]{Produit scientifique}
Calculer $(6\times10^8)(4\times10^{-5})$ en notation scientifique.
\end{exercice}\end{blocLivre}
\begin{blocLivre}{ex-2-7}\begin{exercice}[2]{Quotient scientifique}
Calculer $\dfrac{8{,}4\times10^{-3}}{2{,}1\times10^5}$.
\end{exercice}\end{blocLivre}
\begin{blocLivre}{ex-2-8}\begin{exercice}[2]{Encadrer une racine}
Encadrer $\sqrt{37}$ entre deux dixièmes consécutifs, puis donner son arrondi au dixième.
\end{exercice}\end{blocLivre}
\begin{blocLivre}{ex-2-9}\begin{exercice}[2]{Simplifier exactement}
Simplifier $\sqrt{98}$ et $\sqrt{200}$, puis calculer $\sqrt{98}+\sqrt{200}$.
\end{exercice}\end{blocLivre}
\begin{blocLivre}{ex-2-10}\begin{exercice}[2]{Réfuter une règle}
Réfuter $\sqrt{a+b}=\sqrt a+\sqrt b$ en choisissant deux nombres positifs simples.
\end{exercice}\end{blocLivre}

\clearpage\rubrique{Problèmes}
\begin{blocLivre}{ex-2-11}\begin{exercice}[3]{Une population qui double}
Une culture contient 800 cellules. Leur nombre double toutes les 20 minutes, selon un modèle valable pendant l’expérience. Combien de cellules après 2 heures ? Écrire le résultat en notation scientifique.
\end{exercice}\end{blocLivre}
\begin{blocLivre}{ex-2-12}\begin{exercice}[3]{Une masse minuscule}
Un lot de $4\times10^{12}$ particules a une masse totale de $1{,}8$ g. Dans ce modèle, toutes ont la même masse. Calculer la masse d’une particule en grammes, puis en kilogrammes.
\end{exercice}\end{blocLivre}
\begin{blocLivre}{ex-2-13}\begin{exercice}[3]{Une distance astronomique}
Un signal parcourt $3\times10^5$ km par seconde. Il voyage pendant 8 min 20 s. Quelle distance parcourt-il ? Donner aussi cette distance en mètres.
\end{exercice}\end{blocLivre}
\begin{blocLivre}{ex-2-14}\begin{exercice}[3]{Choisir le côté d’une dalle}
Une dalle carrée a une aire de 2 m$^2$. Donner son côté exact en mètres, puis l’arrondir au millimètre. Peut-on le remplacer par 1,4 m pour calculer une aire exacte ?
\end{exercice}\end{blocLivre}
\begin{blocLivre}{ex-2-15}\begin{exercice}[3]{Comparer deux sauvegardes}
Un fichier de $6{,}4\times10^9$ octets est transmis à $8\times10^6$ octets par seconde, sans variation de débit. Calculer le temps. Un autre système annonce 10 minutes : de quel facteur est-il plus rapide ?
\end{exercice}\end{blocLivre}

\clearpage
